\documentclass[11pt,a4paper]{article}

\usepackage[T1]{fontenc}
\usepackage[utf8]{inputenc}
\usepackage{geometry}
\geometry{margin=1in}
\usepackage{booktabs}
\usepackage{longtable}
\usepackage{array}
\usepackage{multirow}
\usepackage{threeparttable}
\usepackage{tabularx}
\usepackage{hyperref}
\usepackage{fancyhdr}
\usepackage{enumitem}
\usepackage{xcolor}

\hypersetup{
	colorlinks=true,
	linkcolor=black,
	urlcolor=blue,
	citecolor=black,
	pdftitle={NEXUS-ATMS Results Report},
	pdfauthor={NEXUS-ATMS},
}

\pagestyle{fancy}
\fancyhf{}
\lhead{NEXUS-ATMS Results Report}
\rhead{April 2026}
\cfoot{\thepage}

\title{\textbf{NEXUS-ATMS Results Report} \\
\large PDF-Style Submission Draft for Research Paper Audit}
\author{Prepared from project artifacts: evaluation results, benchmark notes, runtime logs}
\date{April 7, 2026}

\begin{document}
\maketitle

\section*{Abstract}
This report summarizes the current NEXUS-ATMS prototype, an AI-centric urban traffic management system that combines SUMO-based reinforcement learning, OpenCV/YOLO vehicle sensing, IoT-style fusion, anomaly detection, and a live FastAPI dashboard. The primary result shows a major reduction in average waiting time and queue length compared with a fixed-timing baseline, while maintaining a live operational dashboard and supporting audit/replay functionality.

\section*{1. Experimental Setup}
\begin{table}[h!]
\centering
\caption{Experimental setup used for the reported results.}
\begin{tabular}{p{0.38\linewidth} p{0.50\linewidth}}
\toprule
\textbf{Item} & \textbf{Value} \\
\midrule
Environment & Standalone \texttt{TrafficEnvironment} with 4-phase NEMA control \\
State dimension & 26 \\
Decision interval & 5 seconds \\
Simulation duration & 720 steps per episode \\
RL training timesteps & 50{,}000 \\
Evaluation episodes & 5 \\
Hardware & AMD Ryzen 5 5600H, NVIDIA RTX 2050 (4.3 GB VRAM) \\
Software & Python 3.13.7, PyTorch 2.6.0+cu124, SB3 2.x \\
\bottomrule
\end{tabular}
\end{table}

\section*{2. Main Traffic Control Results}
\begin{table}[h!]
\centering
\caption{Canonical RL evaluation results from \texttt{results/evaluation\_results.json}.}
\begin{threeparttable}
\begin{tabular}{lccc}
\toprule
\textbf{Metric} & \textbf{Fixed-Timing Baseline} & \textbf{RL Agent} & \textbf{Change} \\
\midrule
Mean reward & $-1149.14 \pm 317.28$ & $-20.66 \pm 2.20$ & $+98.20\%$ \\
Average waiting time (s) & 581.31 & 10.23 & $-98.24\%$ \\
Average queue length (veh) & 26.16 & 2.62 & $-89.97\%$ \\
Throughput (veh/hr) & 510.8 & 453.8 & $-11.16\%$ \\
Episode length & 720 & 720 & -- \\
\bottomrule
\end{tabular}
\begin{tablenotes}
\footnotesize
\item The RL policy strongly reduces delay and queueing.
\item Throughput decreases moderately because the reward function prioritizes congestion relief over raw vehicle count.
\end{tablenotes}
\end{threeparttable}
\end{table}

\section*{3. Training Time}
\begin{table}[h!]
\centering
\caption{Training duration summary from project logs.}
\begin{tabular}{p{0.36\linewidth} p{0.22\linewidth} p{0.20\linewidth} p{0.16\linewidth}}
\toprule
\textbf{Model} & \textbf{Mode} & \textbf{Training Size} & \textbf{Time} \\
\midrule
DQN traffic control policy & SUMO demo mode & 50{,}000 timesteps & 18 min 13 s \\
LSTM traffic predictor & CPU training & 50 epochs & About 2 min \\
\bottomrule
\end{tabular}
\end{table}

\section*{4. Prediction and Anomaly Detection Results}
\begin{table}[h!]
\centering
\caption{Forecasting and anomaly-detection performance.}
\begin{tabular}{p{0.36\linewidth} p{0.24\linewidth} p{0.20\linewidth}}
\toprule
\textbf{Component} & \textbf{Metric} & \textbf{Value} \\
\midrule
LSTM predictor & $R^2$ & 0.6126 \\
LSTM predictor & MAE & 0.0746 \\
IsolationForest & Accuracy & 0.850 \\
Autoencoder & Accuracy & 0.900 \\
Ensemble anomaly detector & Precision & 0.840 \\
Ensemble anomaly detector & Recall & 1.000 \\
Ensemble anomaly detector & F1-score & 0.913 \\
\bottomrule
\end{tabular}
\end{table}

\section*{5. Audit Summary}
\begin{table}[h!]
\centering
\caption{Project audit coverage from simulation to live dashboard.}
\begin{tabular}{p{0.20\linewidth} p{0.25\linewidth} p{0.45\linewidth}}
\toprule
\textbf{Layer} & \textbf{Component} & \textbf{What the paper should say} \\
\midrule
Simulation & SUMO / TraCI & RL traffic signal control on single-intersection and grid scenarios \\
Stand-alone control & First-principles traffic model & Queue/demand/reward modelling without SUMO dependency \\
Vision & OpenCV + YOLOv8 & Vehicle detection, lane-level counting, fallback handling \\
IoT fusion & Sensor simulator + fusion engine & Loop, radar, environmental, pedestrian, and emergency inputs \\
Prediction & LSTM & Multi-step short-horizon traffic forecasting \\
Anomaly detection & Statistical + ML ensemble & Real-time safety-critical alerting \\
RL control & DQN / PPO & Baseline-vs-agent comparison for adaptive signal control \\
Backend & FastAPI + WebSocket & Live system-state streaming, audit logs, replay, and control endpoints \\
Dashboard & Frontend map and KPIs & Operational view with sync indicator, health, incidents, and controls \\
\bottomrule
\end{tabular}
\end{table}

\section*{6. Short Interpretation}
The reported results indicate that the RL controller substantially outperforms fixed timing on congestion-related metrics. The most important gain is waiting time reduction, which dropped from 581.31 seconds to 10.23 seconds. Queue length also fell sharply, from 26.16 vehicles to 2.62 vehicles. The throughput reduction is moderate and reflects the reward design, which prioritizes reduced delay and better traffic balance. In addition to the control results, the project includes traffic forecasting, anomaly detection, and a live dashboard, which makes it a complete AI-based traffic management prototype rather than a single-model experiment.

\section*{7. Conclusion}
The current NEXUS-ATMS implementation already supports end-to-end AI traffic management: simulation, control, forecasting, anomaly detection, live dashboarding, and audit/replay. For paper submission, the canonical table source should be the JSON evaluation file, with training time reported from the DQN run logs. This gives a credible and submission-ready results section for professor review.

\end{document}
